%%
%% part3-customization/ch21-themes.tex
%%
%% Chapter 21: Themes — The default theme, its components,
%% and how to switch themes.
%%

\chapter{تم‌ها}
\label{chap:themes}

\section{چرا این موضوع مهم است؟}
\label{sec:themes-why}

یک تم، مجموعه‌ای از تصمیمات طراحی است که
هویت بصری کتاب را تعیین می‌کند: رنگ‌ها،
فونت‌ها، جلد، و چیدمان. با تغییر تم،
می‌توانید کتاب خود را کاملاً متفاوت جلوه
دهید.

\lr{MatinBook} یک سیستم تم ماژولار ارائه
می‌دهد که به شما امکان می‌دهد تم پیش‌فرض را
استفاده کنید یا تم سفارشی بسازید. در این
فصل، با این سیستم آشنا می‌شوید.

\mnote{تم پیش‌فرض \lr{MatinBook}، \lr{default} نام دارد و تمام قابلیت‌های کلاس را فعال می‌کند.}

\section{تم پیش‌فرض}
\label{sec:themes-default}

تم پیش‌فرض \lr{MatinBook}، \lr{default} است.
این تم از سه فایل تشکیل شده است:

\begin{itemize}
    \item \file{mb-theme-colors.sty}: پالت
    رنگی \lr{CMYK}.
    \item \file{mb-theme-cover.sty}: طراحی
    جلد.
    \item \file{mb-theme-default.sty}:
    بارگذارنده‌ی تم.
\end{itemize}

\mnote{تم \lr{default} از ۲۰ رنگ \lr{CMYK} در ۵ خانواده استفاده می‌کند که مطابق استاندارد نشر آموزشی ایران است.}

\subsection{ساختار فایل‌های تم}
\label{sec:themes-structure}

تم‌ها در پوشه‌ی \file{tex/themes/} قرار
دارند:

\begin{latin}
\begin{minted}{text}
tex/themes/
└── default/
    ├── mb-theme-colors.sty
    ├── mb-theme-cover.sty
    └── mb-theme-default.sty
\end{minted}
\end{latin}

\mnote{هر تم، یک پوشه‌ی جداگانه با سه فایل \lr{.sty} دارد.}

\section{بارگذاری تم}
\label{sec:themes-loading}

\subsection{تم پیش‌فرض}
\label{sec:themes-loading-default}

تم پیش‌فرض به‌طور خودکار در \file{matinbook.cls}
بارگذاری می‌شود:

\begin{latin}
\begin{minted}{latex}
% |\pc{در فایل}|\lr{matinbook.cls}:
\RequirePackage{mb-theme-colors}
...
\RequirePackage{mb-theme-cover}
\end{minted}
\end{latin}

\mnote{ترتیب بارگذاری مهم است: \lr{mb-theme-colors} باید \lr{قبل از} مصرف‌کنندگان (مثل \lr{mb-boxes}) بارگذاری شود.}

\subsection{تم سفارشی}
\label{sec:themes-loading-custom}

برای بارگذاری یک تم سفارشی، فایل
\file{matinbook.cls} را ویرایش کنید:

\begin{latin}
\begin{minted}{latex}
% |\pc{به‌جای}|%
\RequirePackage{mb-theme-colors}
\RequirePackage{mb-theme-cover}

% |\pc{بنویسید}|%
\RequirePackage{tex/themes/mytheme/mb-theme-default}
\end{minted}
\end{latin}

\mnote{مسیر تم سفارشی باید نسبت به پوشه‌ی پروژه باشد.}

\section{ساختار یک تم}
\label{sec:themes-components}

هر تم از سه جزء تشکیل شده است:

\subsection{پالت رنگ}
\label{sec:themes-colors}

فایل \file{mb-theme-colors.sty} رنگ‌های تم
را تعریف می‌کند:

\begin{latin}
\begin{minted}{latex}
\providecolor{matinblue}{cmyk}{0.80, 0.30, 0.00, 0.30}
\providecolor{matingreen}{cmyk}{0.80, 0.00, 0.45, 0.30}
\providecolor{matinorange}{cmyk}{0.00, 0.45, 0.85, 0.10}
\providecolor{matinred}{cmyk}{0.00, 0.65, 0.75, 0.10}
\providecolor{matinpurple}{cmyk}{0.20, 0.60, 0.00, 0.30}
\end{minted}
\end{latin}

\mnote{استفاده از \lr{\textbackslash providecolor} به‌جای \lr{\textbackslash definecolor} امکان بازنویسی رنگ‌ها را در سند فراهم می‌کند.}

\subsection{جلد}
\label{sec:themes-cover}

فایل \file{mb-theme-cover.sty} جلد جلو و
عقب را طراحی می‌کند:

\begin{latin}
\begin{minted}{latex}
\newcommand{\makecover}[4]{%
    \begin{titlepage}\thispagestyle{empty}
        \mb@frontart
        ...
    \end{titlepage}%
}

\newcommand{\makebackcover}[1]{%
    \clearpage
    \thispagestyle{empty}%
    \mb@backart
    ...
}
\end{minted}
\end{latin}

\mnote{جلد با \lr{TikZ} ساخته می‌شود و به حداقل دو بار کامپایل نیاز دارد.}

\subsection{بارگذارنده‌ی تم}
\label{sec:themes-loader}

فایل \file{mb-theme-default.sty} اجزای تم
را بارگذاری می‌کند:

\begin{latin}
\begin{minted}{latex}
\RequirePackage{mb-theme-cover}
\end{minted}
\end{latin}

\mnote{فایل \lr{mb-theme-colors} اینجا بارگذاری نمی‌شود، چون در \lr{matinbook.cls} بارگذاری می‌شود.}

\section{ایجاد تم سفارشی}
\label{sec:themes-creating}

برای ایجاد یک تم سفارشی، این مراحل را
دنبال کنید:

\begin{enumerate}
    \item یک پوشه‌ی جدید در
    \file{tex/themes/} بسازید (مثلاً
    \file{mytheme/}).

    \item سه فایل تم را کپی کنید:

    \begin{latin}
\begin{minted}{bash}
cp tex/themes/default/*.sty tex/themes/mytheme/
\end{minted}
    \end{latin}

    \item فایل‌های تم را ویرایش کنید.

    \item در \file{matinbook.cls}، تم جدید را
    بارگذاری کنید.
\end{enumerate}

\mnote{برای شروع، از تم \lr{default} کپی بگیرید و آن را تغییر دهید. این ساده‌تر از نوشتن تم از صفر است.}

\subsection{تغییر رنگ‌ها}
\label{sec:themes-custom-colors}

ساده‌ترین راه برای تغییر تم، تغییر رنگ‌ها
است. مثلاً برای یک تم با رنگ اصلی سبز:

\begin{latin}
\begin{minted}{latex}
% |\pc{در فایل}|\lr{mb-theme-colors.sty}:
\providecolor{matinblue}{cmyk}{0.80, 0.00, 0.45, 0.30}  % |\pc{سبز}|%
\providecolor{matindarkblue}{cmyk}{0.80, 0.00, 0.45, 0.60}
\providecolor{matinlightblue}{cmyk}{0.10, 0.00, 0.05, 0.05}
\end{minted}
\end{latin}

\mnote{به‌جای تغییر نام رنگ‌ها، فقط مقادیر \lr{CMYK} آن‌ها را تغییر دهید. این کار، سازگاری با کلاس را حفظ می‌کند.}

\subsection{تغییر جلد}
\label{sec:themes-custom-cover}

برای تغییر جلد، فایل \file{mb-theme-cover.sty}
را ویرایش کنید. مثلاً برای تغییر رنگ
پس‌زمینه‌ی جلد:

\begin{latin}
\begin{minted}{latex}
\providecolor{matinnavy}{RGB}{30, 42, 58}
% |\pc{تغییر به قرمز تیره}|%
\providecolor{matinnavy}{RGB}{120, 30, 30}
\end{minted}
\end{latin}

\mnote{برای تغییرات ساده، فقط رنگ‌ها را تغییر دهید. برای طراحی کاملاً جدید، عناصر \lr{TikZ} را جابه‌جا یا حذف کنید.}

\section{چند تم نمونه}
\label{sec:themes-examples}

\subsection{تم \lr{minimal}}
\label{sec:themes-minimal}

یک تم مینیمال که فقط از دو رنگ استفاده
می‌کند:

\begin{latin}
\begin{minted}{latex}
\providecolor{matinblue}{cmyk}{0.00, 0.00, 0.00, 0.80}
\providecolor{matingreen}{cmyk}{0.00, 0.00, 0.00, 0.60}
\providecolor{matinorange}{cmyk}{0.00, 0.00, 0.00, 0.50}
\end{minted}
\end{latin}

\mnote{تم‌های مینیمال برای کتاب‌های آکادمیک و رسمی مناسب‌ترند.}

\subsection{تم \lr{vibrant}}
\label{sec:themes-vibrant}

یک تم پرانرژی با رنگ‌های زنده:

\begin{latin}
\begin{minted}{latex}
\providecolor{matinblue}{cmyk}{1.00, 0.00, 0.00, 0.00}
\providecolor{matingreen}{cmyk}{1.00, 0.00, 1.00, 0.00}
\providecolor{matinorange}{cmyk}{0.00, 0.50, 1.00, 0.00}
\end{minted}
\end{latin}

\mnote{تم‌های پرانرژی برای کتاب‌های آموزشی و مخاطبان جوان مناسب‌ترند.}

\section{خطاهای رایج}
\label{sec:themes-mistakes}

\subsection{عدم بارگذاری \file{mb-theme-colors}}
\label{sec:themes-mistake-colors}

\textbf{مشکل:} اگر \file{mb-theme-colors} را
بارگذاری نکنید، رنگ‌ها تعریف نمی‌شوند و
خطای \lr{Undefined color} می‌دهید.

\textbf{راه‌حل:} مطمئن شوید که
\file{mb-theme-colors} در فاز ۱ \lr{matinbook.cls}
بارگذاری می‌شود.

\subsection{ترتیب اشتباه بارگذاری}
\label{sec:themes-mistake-order}

\textbf{مشکل:} اگر \file{mb-theme-cover} را
قبل از \file{mb-theme-colors} بارگذاری کنید،
خطا می‌دهید.

\textbf{راه‌حل:} ترتیب صحیح:
\file{mb-theme-colors} → \file{mb-theme-cover}.

\subsection{تغییر نام رنگ‌ها}
\label{sec:themes-mistake-rename}

\textbf{مشکل:} اگر نام رنگ‌ها را تغییر دهید،
ماژول‌های مصرف‌کننده (مثل \file{mb-boxes})
خطا می‌دهند.

\textbf{راه‌حل:} نام‌ها را تغییر ندهید؛ فقط
مقادیر \lr{CMYK} را تغییر دهید.

\section{تمرین‌ها}
\label{sec:themes-exercises}

\begin{exercise}
    یک تم سفارشی بسازید که از رنگ‌های
    متفاوتی استفاده کند.
    \label{exc:themes-custom}
\end{exercise}

\begin{exercise}
    رنگ اصلی تم را از آبی به سبز تغییر
    دهید.
    \label{exc:themes-green}
\end{exercise}

\begin{exercise}
    یک تم مینیمال بسازید که فقط از سه
    رنگ خاکستری استفاده کند.
    \label{exc:themes-minimal}
\end{exercise}

\begin{exercise}
    تم خود را در \file{matinbook.cls} بارگذاری
    کنید و کتاب را کامپایل کنید.
    \label{exc:themes-load}
\end{exercise}

\section{خلاصه}
\label{sec:themes-summary}

در این فصل، با تم‌ها در \lr{MatinBook}
آشنا شدیم:

\begin{itemize}
    \item تم پیش‌فرض \lr{default} است و از
    سه فایل تشکیل شده.

    \item هر تم در پوشه‌ی
    \file{tex/themes/} قرار دارد.

    \item اجزای تم: \file{mb-theme-colors}،
    \file{mb-theme-cover}،
    \file{mb-theme-default}.

    \item برای ایجاد تم سفارشی، پوشه‌ی جدید
    بسازید و سه فایل را کپی کنید.

    \item ساده‌ترین راه تغییر تم، تغییر
    رنگ‌ها است.

    \item نام رنگ‌ها را تغییر ندهید؛ فقط
    مقادیر \lr{CMYK} را تغییر دهید.

    \item خطاهای رایج شامل عدم بارگذاری
    \file{mb-theme-colors}، ترتیب اشتباه
    بارگذاری، و تغییر نام رنگ‌ها است.
\end{itemize}

در فصل بعدی (\cref{chap:localization})،
با لوکال و زبان در \lr{MatinBook} آشنا
می‌شوید و یاد می‌گیرید که چگونه بین فارسی
و انگلیسی سوئیچ کنید.